% Paper A draft. Structure follows paper_outline_v3.md; experiments and
% sequencing are in paper_plan.md.
%
% Sections that are settled are written as prose. Everything else sits in an
% \begin{outline} ... \end{outline} block, which prints in a grey frame headed
% "Outline (not yet written)". Numbers that have not been measured on the
% paper's design are written \tbd{..}.
\documentclass[11pt]{article}

\usepackage{amsmath,amssymb,amsthm}
\usepackage{booktabs}
\usepackage{graphicx}
\usepackage[margin=1in]{geometry}
\usepackage[authoryear,round]{natbib}
\usepackage{xcolor}
\usepackage[most]{tcolorbox}
\usepackage{enumitem}
\usepackage{hyperref}

\newtheorem{theorem}{Theorem}
\newtheorem{corollary}{Corollary}
\newtheorem{proposition}{Proposition}
\theoremstyle{definition}
\newtheorem{assumption}{Assumption}
\newtheorem{definition}{Definition}
\theoremstyle{remark}
\newtheorem{remark}{Remark}

\newtcolorbox{outline}{breakable, colback=black!4, colframe=black!35,
  title={Outline (not yet written)}, fonttitle=\small\bfseries,
  left=4pt, right=4pt, top=3pt, bottom=3pt}
\newcommand{\tbd}[1]{\textcolor{red!70!black}{[#1]}}
\setlist{itemsep=2pt, topsep=3pt}

\title{Constraint Learning under Label Noise:\\
A Robust Counterpart that Stays Feasible at the Optimum}
\author{Josephine Situ \and Dimitris Bertsimas}
\date{Draft of \today}

\begin{document}
\maketitle

\begin{abstract}
In constraint learning, a constraint with no known form is replaced by a model
fitted on data and embedded in a mixed-integer optimization problem. When the
training labels are noisy, the optimizer selects the point where the model errs
in its favour, so the prescribed decision often violates the true constraint.
We propose a robust counterpart whose margin $R\|s(x)\|_2$ grows with the
sensitivity of the prediction at $x$ to the training labels. It applies to
ridge regression, trees and forests, and last-layer networks, uses every
training row, and is solved as a mixed-integer second-order cone problem or by
finitely many cuts. We prove that, at any feasible decision, however it was
chosen, the probability of violating the constraint of the model fitted on
noise-free labels is at most $\epsilon$, and give a radius computed by
simulation from the fitted model and the noise level. On a membrane reactor
design problem with a mechanistic ground truth, the nominal decision is
feasible in \tbd{N}\% of runs and the robust counterpart in \tbd{X}\% at a
target of 90\%, at \tbd{Z}\% lower cost than a constant margin, the bootstrap
wrapper and conformal constraint learning at the same feasibility. On
chemotherapy regimen design with real trial data and five learned toxicity
constraints, it certifies a regimen at the 90\% level for \tbd{A}\% of
cohorts, against \tbd{C}\% for the best baseline; every other cohort receives
the regimen with the highest certified level, or is reported as one for which
the data support no safe regimen.
\end{abstract}

% ---------------------------------------------------------------------------
\section{Introduction}

\begin{outline}
\begin{itemize}
  \item Open with the failure, in numbers. On the reactor under the cost
    vectors of \citet{ovalle2025cmicl}, the nominal decision is feasible under
    the ODE in only \tbd{0.18 to 0.25} of runs with their network and
    \tbd{0.66} with cubic ridge (probes, 100 noise redraws; final numbers
    depend on the lead model).
  \item The cleanest illustration of selection is C-MICL itself: at random
    points inside its feasible region the truth is feasible \tbd{0.93 to 0.99}
    of the time for their network, forest and boosting surrogates, and at the
    optimizer's choice \tbd{0.59 to 0.72} (Section~\ref{sec:experiments},
    preliminary results).
  \item Why it happens: the optimizer searches for the point where the model
    is most optimistic. A bound that holds at a fixed $x$ does not hold at
    $x^\star$. This is the optimizer's curse \citep{smith2006optimizer}
    applied to feasibility rather than to the objective;
    \citet{cremer2019datadriven} report a 99.9\%-accurate learned security
    rule that is 0\% accurate at the optimized operating points.
  \item Why existing safeguards fall short. A constant margin pays the same
    everywhere and has no rule for its size. The wrapper
    \citep{maragno2025mixed} has no link between its fraction and a violation
    level. C-MICL \citep{ovalle2025cmicl} varies its margin with $x$ through a
    second, learned model, whose own errors the optimizer can exploit; it also
    holds out calibration rows and doubles the embedded model. Its guarantee
    holds for an exchangeable test point and reaches the optimum only through
    a conditional-independence assumption that the authors say may fail.
  \item What we do: a margin that varies with the decision and is computed,
    not learned.
  \item What a practitioner gets: a radius set from the target violation
    level, with no ground truth, no calibration split and no tuning run.
  \item Contributions (four): the counterpart; exact solution methods; the
    guarantee at $x^\star$ with a simulated radius; experiments \tbd{to run}.
  \item Literature, one page (checked 2026-11-14; notes in
    \texttt{paper\_a\_literature.md}).
    \begin{itemize}
      \item Constraint learning \citep{fajemisin2024optimization,
        maragno2025mixed}; conformal \citep{ovalle2025cmicl}. Other
        uncertainty-aware constraint learning, none modelling the error of
        the fitted model under label noise: conditional-quantile models
        \citep{alcantara2023datadriven} hedge the outcome's own spread, not
        the estimate's error; \citet{cremer2019datadriven} tune per-split
        margins on embedded trees against a simulator; DCL
        \citep{oztekin2026dcl} localizes the surrogate and leaves feasibility
        certification open.
      \item Robust optimization with statistical guarantees
        \citep{bertsimas2022robust, bertsimas2018datadriven,
        hong2021learning}. For a linear model the label ball is an ellipsoid
        on the coefficients, so with the closed-form radius our counterpart is
        the robust counterpart of a confidence ellipsoid
        \citep{goldfarb2003robust}; RO on surrogate parameters
        \citep{hullen2020managing}. Cutting-set methods
        \citep{mutapcic2009cutting, bertsimas2016reformulation}.
      \item Simultaneous inference. Bands restricted to a region of the
        covariates \citep{wynn1971simultaneous, naiman1987simultaneous}, with
        the critical constant simulated by optimizing over the region on each
        draw \citep{liu2005simulation, liu2010simultaneous}: this is our
        simulated radius for a linear model over a box. Sup-$t$ bands
        \citep{montielolea2019simultaneous}; Scheff\'e
        \citep{scheffe1953method}. Validity at a data-selected point from a
        simultaneous band is the projection argument of post-selection and
        winner's-curse inference \citep{berk2013valid, andrews2024inference}.
      \item Uniform bounds in safe learning. ``Mean $+\,\beta\cdot$sd $\le b$''
        with $\beta$ valid over the whole domain is the safe set of safe
        Bayesian optimization \citep{srinivas2012information, sui2015safe,
        berkenkamp2023bayesian} and of safe linear bandits
        \citep{amani2019linear}. \citet{fiedler2021practical} bound the noise
        term of kernel ridge regression by a multiple of the same
        $\|s(x)\|_2$, uniformly over the domain: the closest precedent for
        the margin, though for learning-based robust control (a learned
        dynamics model inside a robust controller), not constraint
        learning. GP
        chance constraints in global optimization
        \citep{schweidtmann2021deterministic, wiebe2022robust} have the same
        algebraic form with a Bayesian, pointwise probability; in mixed-integer
        form over trees the multiplier is a tuning constant
        \citep{thebelt2022tree}.
      \item What is new, stated against these: the target is the clean-label
        model, so no bound on the true function's norm enters the radius (and
        the claim is weaker, which we say); the region is the mixed-integer
        feasible set and each simulated supremum is a global MIP solved to a
        bound; one label-space form covers ridge, last-layer networks and
        tree ensembles; and the band is imposed as a constraint, so the
        optimizer's choice inherits it. Not new, and not claimed: restricting
        a band to a region, simulating its constant by optimization, the
        ellipsoidal counterpart for linear models, or validity at a selected
        point by simultaneity.
    \end{itemize}
\end{itemize}
\end{outline}

% ---------------------------------------------------------------------------
\section{Problem and method}\label{sec:method}

\subsection{Constraint learning}
A decision $x \in \mathcal{X} \subseteq \mathbb{R}^p$ minimizes a known cost
subject to known constraints, which define $\mathcal{X}$, and one constraint
$h^\star(x) \le b$ whose function $h^\star$ is not known. Constraint learning
\citep{fajemisin2024optimization, maragno2025mixed} replaces $h^\star$ by a
model $f(\cdot\,;\theta)$ fitted on data $(x_i, y_i)_{i=1}^n$ and embeds the
fitted model in the optimization problem:
\begin{equation}\label{eq:nominal}
  \min_{x \in \mathcal{X}} \; c^\top x
  \quad\text{s.t.}\quad f(x;\hat\theta) \le b .
\end{equation}
We call \eqref{eq:nominal} the nominal problem. The set $\mathcal{X}$ includes
any trust region that keeps $x$ near the training inputs.

\subsection{Label noise}
The training inputs $x_i$ are observed exactly. The labels are observed with
noise: $y = y^\star + \varepsilon$, where $y^\star_i = h^\star(x_i)$. We write
$\theta(y)$ for the model fitted on labels $y$, $\hat\theta = \theta(y)$ for
the model we have, and $\theta^\star = \theta(y^\star)$ for the model the same
procedure would return on noise-free labels. We call $f(\cdot\,;\theta^\star)$
the clean-label model.

\begin{outline}
One sentence per case study on where the noise comes from: measurement noise
in the reactor, sampling error of a trial arm's toxicity rate, survey noise in
palatability.
\end{outline}

\subsection{The sensitivity vector}
Many regression models are affine in their training labels once their
structure is fixed: there is a vector $s(x) \in \mathbb{R}^n$ and a scalar
$c(x)$, neither depending on $y$, with
\begin{equation}\label{eq:affine}
  f(x;\theta(y)) = s(x)^\top y + c(x) .
\end{equation}
The entry $s_i(x)$ is the weight that training label $i$ carries in the
prediction at $x$. We call $s(x)$ the sensitivity vector.

\begin{itemize}
  \item \emph{Ridge regression} on features $\phi(x)$, with training matrix
    $Z$ and penalty $\lambda$: $s(x) = Z(Z^\top Z + \lambda I)^{-1}\phi(x)$.
    Its squared norm is, up to the penalty, the leverage of $x$.
  \item \emph{Last-layer network}: the hidden layers are frozen and the output
    layer is fitted by ridge regression, so $\phi(x)$ is the last hidden
    layer's activation and $s(x)$ is as above.
  \item \emph{Regression tree} with fixed splits: $s_i(x) = 1/n_\ell$ if row
    $i$ falls in the leaf $\ell$ that contains $x$ and $0$ otherwise, where
    $n_\ell$ is the number of training rows in that leaf. So
    $\|s(x)\|_2 = 1/\sqrt{n_\ell}$.
  \item \emph{Forest} with fixed splits: the average of its trees' vectors,
    which are the forest weights of \citet{meinshausen2006quantile}.
\end{itemize}

\subsection{The robust counterpart}
Suppose the noise vector lies in the ball $\{\delta : \|\delta\|_2 \le R\}$.
The clean labels are then $y - \delta$ for some $\delta$ in the ball, and by
\eqref{eq:affine} the clean-label prediction is $f(x;\hat\theta) -
s(x)^\top\delta$. Requiring the constraint for every such $\delta$ gives
\begin{equation}\label{eq:rc-method}
  f(x;\hat\theta) + R\,\|s(x)\|_2 \le b ,
\end{equation}
since $\max_{\|\delta\|_2 \le R} \, -s(x)^\top\delta = R\|s(x)\|_2$. This is
the exact robust counterpart; nothing is approximated. The robust problem is
\eqref{eq:nominal} with \eqref{eq:rc-method} in place of the learned
constraint. For ridge regression the coefficients are linear in the labels,
so the ball on the labels is an ellipsoid on the coefficients and
\eqref{eq:rc-method} is the familiar counterpart of ellipsoidal parameter
uncertainty \citep{goldfarb2003robust, bertsimas2022robust}. What the
label-space form adds is that the same constraint applies unchanged to trees,
forests and last-layer networks.

The margin $R\|s(x)\|_2$ depends on the decision. For a tree it is
$R/\sqrt{n_\ell}$, a multiple of the standard error of the leaf mean. For ridge
regression it is a multiple of the square root of the leverage. In both cases
it is large where few training labels determine the prediction, which is where
the optimizer is extrapolating, and small where many do. It is computed from
the fitted model and the training inputs; no second model is trained.

\begin{outline}
\begin{itemize}
  \item Weighted norm when the noise levels $\sigma_i$ are known and unequal
    (trial arms of different sizes).
  \item Figure 1: the reactor's feasible region under the nominal model, a
    constant margin and the robust counterpart, with the true boundary.
  \item Choosing $R$, the recipe: (0) fix the model before looking at the
    labels; (1) estimate $\sigma$ from out-of-fold residuals; (2) simulate the
    radius for the target $\epsilon$ (Section~\ref{sec:guarantee}); (3) solve.
    No ground truth and no calibration rows are used.
\end{itemize}
\end{outline}

% ---------------------------------------------------------------------------
\section{Guarantee at the optimum}\label{sec:guarantee}

A confidence bound on $f(x;\theta^\star)$ at a fixed $x$ does not hold at the
optimizer's $x^\star$, because $x^\star$ is chosen using the same noisy labels
\citep{smith2006optimizer, andrews2024inference}. This section bounds the
violation probability at any feasible point, however it was chosen.

% Paper A: selection-robust violation guarantee.
% Included by paper.tex (% Paper A: selection-robust violation guarantee.
% Included by paper.tex (% Paper A: selection-robust violation guarantee.
% Included by paper.tex (\input{paper_a_guarantee}) in Section 3. Uses the
% assumption, definition, theorem, corollary and remark environments and the
% \status macro defined there. Citation keys are those of paper_a.bib.
%
% Proofs other than Theorem 1's are in paper_a_proofs.tex (the appendix).
% Results held until probe P1 decides the chemotherapy design (joint radius,
% certified level, simulated level) are in paper_a_guarantee_held.tex, which
% is not \input.
%
% \status tags. JS checked by hand the version committed in 40acea5: Theorem 1,
% Corollaries 1-3 and the closing remark, for Sigma = sigma^2 I. The general
% Sigma, the remark on the choice of norm and everything from "The simulated
% radius" on came later and are not checked.

\begin{outline}
Tags: JS checked by hand the $\Sigma = \sigma^2 I$ version of
Theorem~\ref{thm:selection}, Corollaries~\ref{cor:finite}--\ref{cor:multi}
and the closing remark (commit \texttt{40acea5}). The general $\Sigma$, the
remark on the choice of norm and everything after the closing remark were
added later and are not checked.
\end{outline}

\paragraph{Setting.}
We observe training data $(x_i, y_i)_{i=1}^n$ with $y = y^\star + \varepsilon$,
where $y^\star_i = h^\star(x_i)$ are the noise-free labels. The learned
constraint is $f(x;\theta(y)) \le b$ for $x \in \mathcal{X}$, where $\theta(y)$
denotes the model fitted on labels $y$.

\begin{assumption}[Label noise]\label{ass:noise}
\status{checked by JS for $\Sigma = \sigma^2 I$; general $\Sigma$ not checked}
The noise is $\varepsilon = \Sigma^{1/2}\xi$, where $\Sigma \succeq 0$ is known
and $\xi_1,\dots,\xi_n$ are independent, zero-mean and sub-Gaussian with
variance proxy $1$. We write $\|v\|_\Sigma = \|\Sigma^{1/2}v\|_2$.
\end{assumption}

With $\Sigma = \sigma^2 I$ the noise terms are independent with variance proxy
$\sigma^2$ and $\|v\|_\Sigma = \sigma\|v\|_2$. A diagonal $\Sigma$ allows
unequal noise levels, and a full one correlated noise.

\begin{assumption}[Affine predictor]\label{ass:affine}
For all $x \in \mathcal{X}$, $f(x;\theta(y)) = s(x)^\top y + c(x)$, where
$s(x) \in \mathbb{R}^n$ and $c(x)$ do not depend on $\varepsilon$.
\end{assumption}

\begin{definition}[Clean-label model]\label{def:clean}
Let $\hat\theta = \theta(y)$ and $\theta^\star = \theta(y^\star)$. The
clean-label model is $f(\cdot\,;\theta^\star)$.
\end{definition}

Since $s(x)$ and $c(x)$ do not depend on $\varepsilon$,
Assumption~\ref{ass:affine} applies to $y^\star = y - \varepsilon$ as well as to
$y$, and subtracting the two gives, for all $x \in \mathcal{X}$,
\begin{equation}\label{eq:identity}
  f(x;\theta^\star) = f(x;\hat\theta) - s(x)^\top \varepsilon .
\end{equation}

For a radius $R > 0$ fixed independently of $\varepsilon$, the robust
counterpart of the learned constraint is
\begin{equation}\label{eq:rc}
  f(x;\hat\theta) + R\,\|s(x)\|_\Sigma \le b .
\end{equation}
The radius $R$ is in units of the noise's scale: with $\Sigma = \sigma^2 I$ the
margin is $R\sigma\|s(x)\|_2$.

\begin{theorem}\label{thm:selection}
\status{checked by JS for $\Sigma = \sigma^2 I$; general $\Sigma$ not checked}
Suppose Assumptions~\ref{ass:noise} and~\ref{ass:affine} hold, and let $x^\star$
be any solution satisfying \eqref{eq:rc}, possibly chosen as a function of the
data. Then
\[
  \mathbb{P}\bigl(f(x^\star;\theta^\star) > b\bigr) \le \mathbb{P}(W > R),
  \qquad
  W = \sup_{x \in \mathcal{X},\, \|s(x)\|_\Sigma \neq 0}
      \frac{-s(x)^\top \varepsilon}{\|s(x)\|_\Sigma} .
\]
\end{theorem}

\begin{proof}
Suppose $f(x^\star;\theta^\star) > b$. By \eqref{eq:identity} and
\eqref{eq:rc},
\[
  -s(x^\star)^\top \varepsilon
  = f(x^\star;\theta^\star) - f(x^\star;\hat\theta)
  > b - f(x^\star;\hat\theta)
  \ge R\,\|s(x^\star)\|_\Sigma .
\]
The left side equals $-(\Sigma^{1/2}s(x^\star))^\top\xi$ and is nonzero, so
$\|s(x^\star)\|_\Sigma \neq 0$ and $W > R$.
\end{proof}

\begin{remark}[Choice of norm]\label{rem:norm}
\status{not checked: added after JS's check}
The proof uses only the inequality
$-s(x^\star)^\top\varepsilon > R\,\|s(x^\star)\|_\Sigma$, so
Theorem~\ref{thm:selection} holds verbatim for any norm $N$ on $\mathbb{R}^n$
in place of $\|\cdot\|_\Sigma$, both in \eqref{eq:rc} and in $W$. The norm
enters only through the distribution of $W$. Under $\|\cdot\|_\Sigma$ the
standardized error at $x$ is
$-s(x)^\top\varepsilon/\|s(x)\|_\Sigma = -u(x)^\top\xi$ with $u(x)$ a unit
vector, so it has the same tail at every $x$; Corollaries~\ref{cor:finite}
and~\ref{cor:subspace} use exactly this. Under another norm the standardized
error at $x$ is that one multiplied by
$\rho(x) = \|s(x)\|_\Sigma / N(s(x))$. Bounding the corresponding $W_N$
through $W$, as $W_N \le \bar\rho\max(W,0)$ with
$\bar\rho = \sup_{x\in\mathcal{X}} \rho(x)$, turns a radius $R$ from these
corollaries into $\bar\rho R$, and the margin
$\bar\rho R\,N(s(x)) \ge R\,\|s(x)\|_\Sigma$ is at least as large at every
$x$. For example, noise bounded as $|\varepsilon_i| \le c$ gives the box
counterpart $f(x;\hat\theta) + c\,\|s(x)\|_1 \le b$, which holds with
probability one; but $\|s(x)\|_1$ can be up to $\sqrt{n}$ times
$\|s(x)\|_2$.
\end{remark}

\begin{corollary}[Tree ensembles]\label{cor:finite}
\status{checked by JS for $\Sigma = \sigma^2 I$; general $\Sigma$ not checked}
If $s(x)$ takes at most $K$ distinct values on $\mathcal{X}$, then
$R = \sqrt{2\ln(K/\epsilon)}$ guarantees
$\mathbb{P}\bigl(f(x^\star;\theta^\star) > b\bigr) \le \epsilon$.
\end{corollary}

\begin{corollary}[Linear models]\label{cor:subspace}
\status{checked by JS for $\Sigma = \sigma^2 I$; general $\Sigma$ not checked}
Suppose $s(x)$ lies in a subspace $S \subseteq \mathbb{R}^n$ of dimension $d$
for all $x \in \mathcal{X}$, and $\xi \sim N(0, I)$, that is,
$\varepsilon \sim N(0,\Sigma)$. Then $R = \sqrt{\chi^2_{d,1-\epsilon}}$
guarantees $\mathbb{P}\bigl(f(x^\star;\theta^\star) > b\bigr) \le \epsilon$,
where $\chi^2_{d,1-\epsilon}$ is the $(1-\epsilon)$-quantile of the chi-square
distribution with $d$ degrees of freedom.
\end{corollary}

\begin{corollary}[Multiple learned constraints]\label{cor:multi}
\status{checked by JS for $\Sigma = \sigma^2 I$; general $\Sigma$ not checked}
Consider $m$ learned constraints, each satisfying
Assumptions~\ref{ass:noise} and~\ref{ass:affine} with its own $s_j$,
$\Sigma_j$ and radius $R_j$. If each $R_j$ is chosen as in
Corollary~\ref{cor:finite} or~\ref{cor:subspace} at level $\epsilon/m$, then
the probability that $x^\star$ violates at least one clean-label constraint is
at most $\epsilon$. No assumption is required on the dependence of the noise
across constraints.
\end{corollary}

\begin{remark}
\status{checked by JS for $\Sigma = \sigma^2 I$}
(i) For ridge regression, $s(x) = Z(Z^\top Z + \lambda I)^{-1}\phi(x)$, so $S$
is the column space of the training matrix $Z$ and $d = \mathrm{rank}(Z)$.
(ii) A lower-bound constraint $f(x;\theta) \ge b$ is handled symmetrically,
with robust counterpart $f(x;\hat\theta) - R\|s(x)\|_\Sigma \ge b$ and the same
radii.
(iii) The guarantee concerns the propagation of label noise only; the
approximation error $h^\star - f(\cdot\,;\theta^\star)$ is not covered.
(iv) For models whose structure is selected from data (active sets, splits,
hidden layers), Assumption~\ref{ass:affine} requires the structure to be
chosen independently of $\varepsilon$, for instance by sample splitting.
\end{remark}

\paragraph{The simulated radius.}
Corollaries~\ref{cor:finite} and~\ref{cor:subspace} bound $W$ without using
$\mathcal{X}$: they protect against every direction among the vectors
$\Sigma^{1/2}s(x)$, whether or not a decision in $\mathcal{X}$ reaches it. When
the distribution of $\xi$ is known, the radius can instead be read off
simulated copies of $W$, and an order statistic of the copies is valid with no
sampling-error term. Write $W = w(\xi)$, where
\[
  w(z) = \sup_{x \in \mathcal{X},\, \|s(x)\|_\Sigma \neq 0}
    \frac{-(\Sigma^{1/2}s(x))^\top z}{\|s(x)\|_\Sigma},
  \qquad z \in \mathbb{R}^n ,
\]
assumed measurable. Under Assumption~\ref{ass:affine}, $s(\cdot)$ does not
depend on $\varepsilon$, so $w$ is a fixed function, known once the model's
structure is. Evaluating $w(z)$ is a maximization over $\mathcal{X}$, a
mixed-integer program when $\mathcal{X}$ is mixed-integer.

\begin{proposition}[Simulated radius]\label{prop:simulated}
\status{not checked: drafted 2026-11-15}
Suppose Assumptions~\ref{ass:noise} and~\ref{ass:affine} hold and the
distribution of $\xi$ is known. Draw $\xi^{(1)},\dots,\xi^{(N)}$ from it,
independently of each other and of $\varepsilon$, and let
$\overline W_1,\dots,\overline W_N$ be any numbers with
$\overline W_i \ge w(\xi^{(i)})$. For $\epsilon \in (0,1)$ and
$N \ge 1/\epsilon - 1$, let $k = \lceil (N+1)(1-\epsilon) \rceil$ and let
$\hat R$ be the $k$-th smallest of $\overline W_1,\dots,\overline W_N$. Then
any $x^\star$ satisfying \eqref{eq:rc} at $R = \hat R$, possibly chosen as a
function of the data and of the simulation, satisfies
$\mathbb{P}\bigl(f(x^\star;\theta^\star) > b\bigr) \le \epsilon$, where the
probability is over the noise and the simulation jointly.
\end{proposition}

\begin{remark}[Computing the radius]\label{rem:simulated}
\status{not checked: drafted 2026-11-15}
(i) Each $\overline W_i$ must be an upper bound on $w(\xi^{(i)})$. The value
of a feasible point of the maximization, or the largest standardized error
over a sample of decisions, is a lower bound, and an order statistic of lower
bounds can fall below the $(1-\epsilon)$-quantile of $W$, which makes the
radius too small. A solver's dual bound is an upper bound at any time limit,
so stopping a maximization early makes the radius larger, never invalid.
(ii) The proposition needs the distribution of $\xi$, not only a variance
proxy: the argument requires the simulated copies to have the law of the
true $W$. The experiments use $\xi \sim N(0, I)$. An estimated $\Sigma$ is
not covered yet. \tbd{T2}
\end{remark}

Corollary~\ref{cor:multi} holds with unequal levels $\epsilon_j$ summing to
$\epsilon$, and with any radius valid at level $\epsilon_j$ for constraint $j$
alone, including the simulated radius of
Proposition~\ref{prop:simulated}: its proof uses only
$\mathbb{P}(W_j > R_j) \le \epsilon_j$. \status{not checked: drafted
2026-11-15}
) in Section 3. Uses the
% assumption, definition, theorem, corollary and remark environments and the
% \status macro defined there. Citation keys are those of paper_a.bib.
%
% Proofs other than Theorem 1's are in paper_a_proofs.tex (the appendix).
% Results held until probe P1 decides the chemotherapy design (joint radius,
% certified level, simulated level) are in paper_a_guarantee_held.tex, which
% is not \input.
%
% \status tags. JS checked by hand the version committed in 40acea5: Theorem 1,
% Corollaries 1-3 and the closing remark, for Sigma = sigma^2 I. The general
% Sigma, the remark on the choice of norm and everything from "The simulated
% radius" on came later and are not checked.

\begin{outline}
Tags: JS checked by hand the $\Sigma = \sigma^2 I$ version of
Theorem~\ref{thm:selection}, Corollaries~\ref{cor:finite}--\ref{cor:multi}
and the closing remark (commit \texttt{40acea5}). The general $\Sigma$, the
remark on the choice of norm and everything after the closing remark were
added later and are not checked.
\end{outline}

\paragraph{Setting.}
We observe training data $(x_i, y_i)_{i=1}^n$ with $y = y^\star + \varepsilon$,
where $y^\star_i = h^\star(x_i)$ are the noise-free labels. The learned
constraint is $f(x;\theta(y)) \le b$ for $x \in \mathcal{X}$, where $\theta(y)$
denotes the model fitted on labels $y$.

\begin{assumption}[Label noise]\label{ass:noise}
\status{checked by JS for $\Sigma = \sigma^2 I$; general $\Sigma$ not checked}
The noise is $\varepsilon = \Sigma^{1/2}\xi$, where $\Sigma \succeq 0$ is known
and $\xi_1,\dots,\xi_n$ are independent, zero-mean and sub-Gaussian with
variance proxy $1$. We write $\|v\|_\Sigma = \|\Sigma^{1/2}v\|_2$.
\end{assumption}

With $\Sigma = \sigma^2 I$ the noise terms are independent with variance proxy
$\sigma^2$ and $\|v\|_\Sigma = \sigma\|v\|_2$. A diagonal $\Sigma$ allows
unequal noise levels, and a full one correlated noise.

\begin{assumption}[Affine predictor]\label{ass:affine}
For all $x \in \mathcal{X}$, $f(x;\theta(y)) = s(x)^\top y + c(x)$, where
$s(x) \in \mathbb{R}^n$ and $c(x)$ do not depend on $\varepsilon$.
\end{assumption}

\begin{definition}[Clean-label model]\label{def:clean}
Let $\hat\theta = \theta(y)$ and $\theta^\star = \theta(y^\star)$. The
clean-label model is $f(\cdot\,;\theta^\star)$.
\end{definition}

Since $s(x)$ and $c(x)$ do not depend on $\varepsilon$,
Assumption~\ref{ass:affine} applies to $y^\star = y - \varepsilon$ as well as to
$y$, and subtracting the two gives, for all $x \in \mathcal{X}$,
\begin{equation}\label{eq:identity}
  f(x;\theta^\star) = f(x;\hat\theta) - s(x)^\top \varepsilon .
\end{equation}

For a radius $R > 0$ fixed independently of $\varepsilon$, the robust
counterpart of the learned constraint is
\begin{equation}\label{eq:rc}
  f(x;\hat\theta) + R\,\|s(x)\|_\Sigma \le b .
\end{equation}
The radius $R$ is in units of the noise's scale: with $\Sigma = \sigma^2 I$ the
margin is $R\sigma\|s(x)\|_2$.

\begin{theorem}\label{thm:selection}
\status{checked by JS for $\Sigma = \sigma^2 I$; general $\Sigma$ not checked}
Suppose Assumptions~\ref{ass:noise} and~\ref{ass:affine} hold, and let $x^\star$
be any solution satisfying \eqref{eq:rc}, possibly chosen as a function of the
data. Then
\[
  \mathbb{P}\bigl(f(x^\star;\theta^\star) > b\bigr) \le \mathbb{P}(W > R),
  \qquad
  W = \sup_{x \in \mathcal{X},\, \|s(x)\|_\Sigma \neq 0}
      \frac{-s(x)^\top \varepsilon}{\|s(x)\|_\Sigma} .
\]
\end{theorem}

\begin{proof}
Suppose $f(x^\star;\theta^\star) > b$. By \eqref{eq:identity} and
\eqref{eq:rc},
\[
  -s(x^\star)^\top \varepsilon
  = f(x^\star;\theta^\star) - f(x^\star;\hat\theta)
  > b - f(x^\star;\hat\theta)
  \ge R\,\|s(x^\star)\|_\Sigma .
\]
The left side equals $-(\Sigma^{1/2}s(x^\star))^\top\xi$ and is nonzero, so
$\|s(x^\star)\|_\Sigma \neq 0$ and $W > R$.
\end{proof}

\begin{remark}[Choice of norm]\label{rem:norm}
\status{not checked: added after JS's check}
The proof uses only the inequality
$-s(x^\star)^\top\varepsilon > R\,\|s(x^\star)\|_\Sigma$, so
Theorem~\ref{thm:selection} holds verbatim for any norm $N$ on $\mathbb{R}^n$
in place of $\|\cdot\|_\Sigma$, both in \eqref{eq:rc} and in $W$. The norm
enters only through the distribution of $W$. Under $\|\cdot\|_\Sigma$ the
standardized error at $x$ is
$-s(x)^\top\varepsilon/\|s(x)\|_\Sigma = -u(x)^\top\xi$ with $u(x)$ a unit
vector, so it has the same tail at every $x$; Corollaries~\ref{cor:finite}
and~\ref{cor:subspace} use exactly this. Under another norm the standardized
error at $x$ is that one multiplied by
$\rho(x) = \|s(x)\|_\Sigma / N(s(x))$. Bounding the corresponding $W_N$
through $W$, as $W_N \le \bar\rho\max(W,0)$ with
$\bar\rho = \sup_{x\in\mathcal{X}} \rho(x)$, turns a radius $R$ from these
corollaries into $\bar\rho R$, and the margin
$\bar\rho R\,N(s(x)) \ge R\,\|s(x)\|_\Sigma$ is at least as large at every
$x$. For example, noise bounded as $|\varepsilon_i| \le c$ gives the box
counterpart $f(x;\hat\theta) + c\,\|s(x)\|_1 \le b$, which holds with
probability one; but $\|s(x)\|_1$ can be up to $\sqrt{n}$ times
$\|s(x)\|_2$.
\end{remark}

\begin{corollary}[Tree ensembles]\label{cor:finite}
\status{checked by JS for $\Sigma = \sigma^2 I$; general $\Sigma$ not checked}
If $s(x)$ takes at most $K$ distinct values on $\mathcal{X}$, then
$R = \sqrt{2\ln(K/\epsilon)}$ guarantees
$\mathbb{P}\bigl(f(x^\star;\theta^\star) > b\bigr) \le \epsilon$.
\end{corollary}

\begin{corollary}[Linear models]\label{cor:subspace}
\status{checked by JS for $\Sigma = \sigma^2 I$; general $\Sigma$ not checked}
Suppose $s(x)$ lies in a subspace $S \subseteq \mathbb{R}^n$ of dimension $d$
for all $x \in \mathcal{X}$, and $\xi \sim N(0, I)$, that is,
$\varepsilon \sim N(0,\Sigma)$. Then $R = \sqrt{\chi^2_{d,1-\epsilon}}$
guarantees $\mathbb{P}\bigl(f(x^\star;\theta^\star) > b\bigr) \le \epsilon$,
where $\chi^2_{d,1-\epsilon}$ is the $(1-\epsilon)$-quantile of the chi-square
distribution with $d$ degrees of freedom.
\end{corollary}

\begin{corollary}[Multiple learned constraints]\label{cor:multi}
\status{checked by JS for $\Sigma = \sigma^2 I$; general $\Sigma$ not checked}
Consider $m$ learned constraints, each satisfying
Assumptions~\ref{ass:noise} and~\ref{ass:affine} with its own $s_j$,
$\Sigma_j$ and radius $R_j$. If each $R_j$ is chosen as in
Corollary~\ref{cor:finite} or~\ref{cor:subspace} at level $\epsilon/m$, then
the probability that $x^\star$ violates at least one clean-label constraint is
at most $\epsilon$. No assumption is required on the dependence of the noise
across constraints.
\end{corollary}

\begin{remark}
\status{checked by JS for $\Sigma = \sigma^2 I$}
(i) For ridge regression, $s(x) = Z(Z^\top Z + \lambda I)^{-1}\phi(x)$, so $S$
is the column space of the training matrix $Z$ and $d = \mathrm{rank}(Z)$.
(ii) A lower-bound constraint $f(x;\theta) \ge b$ is handled symmetrically,
with robust counterpart $f(x;\hat\theta) - R\|s(x)\|_\Sigma \ge b$ and the same
radii.
(iii) The guarantee concerns the propagation of label noise only; the
approximation error $h^\star - f(\cdot\,;\theta^\star)$ is not covered.
(iv) For models whose structure is selected from data (active sets, splits,
hidden layers), Assumption~\ref{ass:affine} requires the structure to be
chosen independently of $\varepsilon$, for instance by sample splitting.
\end{remark}

\paragraph{The simulated radius.}
Corollaries~\ref{cor:finite} and~\ref{cor:subspace} bound $W$ without using
$\mathcal{X}$: they protect against every direction among the vectors
$\Sigma^{1/2}s(x)$, whether or not a decision in $\mathcal{X}$ reaches it. When
the distribution of $\xi$ is known, the radius can instead be read off
simulated copies of $W$, and an order statistic of the copies is valid with no
sampling-error term. Write $W = w(\xi)$, where
\[
  w(z) = \sup_{x \in \mathcal{X},\, \|s(x)\|_\Sigma \neq 0}
    \frac{-(\Sigma^{1/2}s(x))^\top z}{\|s(x)\|_\Sigma},
  \qquad z \in \mathbb{R}^n ,
\]
assumed measurable. Under Assumption~\ref{ass:affine}, $s(\cdot)$ does not
depend on $\varepsilon$, so $w$ is a fixed function, known once the model's
structure is. Evaluating $w(z)$ is a maximization over $\mathcal{X}$, a
mixed-integer program when $\mathcal{X}$ is mixed-integer.

\begin{proposition}[Simulated radius]\label{prop:simulated}
\status{not checked: drafted 2026-11-15}
Suppose Assumptions~\ref{ass:noise} and~\ref{ass:affine} hold and the
distribution of $\xi$ is known. Draw $\xi^{(1)},\dots,\xi^{(N)}$ from it,
independently of each other and of $\varepsilon$, and let
$\overline W_1,\dots,\overline W_N$ be any numbers with
$\overline W_i \ge w(\xi^{(i)})$. For $\epsilon \in (0,1)$ and
$N \ge 1/\epsilon - 1$, let $k = \lceil (N+1)(1-\epsilon) \rceil$ and let
$\hat R$ be the $k$-th smallest of $\overline W_1,\dots,\overline W_N$. Then
any $x^\star$ satisfying \eqref{eq:rc} at $R = \hat R$, possibly chosen as a
function of the data and of the simulation, satisfies
$\mathbb{P}\bigl(f(x^\star;\theta^\star) > b\bigr) \le \epsilon$, where the
probability is over the noise and the simulation jointly.
\end{proposition}

\begin{remark}[Computing the radius]\label{rem:simulated}
\status{not checked: drafted 2026-11-15}
(i) Each $\overline W_i$ must be an upper bound on $w(\xi^{(i)})$. The value
of a feasible point of the maximization, or the largest standardized error
over a sample of decisions, is a lower bound, and an order statistic of lower
bounds can fall below the $(1-\epsilon)$-quantile of $W$, which makes the
radius too small. A solver's dual bound is an upper bound at any time limit,
so stopping a maximization early makes the radius larger, never invalid.
(ii) The proposition needs the distribution of $\xi$, not only a variance
proxy: the argument requires the simulated copies to have the law of the
true $W$. The experiments use $\xi \sim N(0, I)$. An estimated $\Sigma$ is
not covered yet. \tbd{T2}
\end{remark}

Corollary~\ref{cor:multi} holds with unequal levels $\epsilon_j$ summing to
$\epsilon$, and with any radius valid at level $\epsilon_j$ for constraint $j$
alone, including the simulated radius of
Proposition~\ref{prop:simulated}: its proof uses only
$\mathbb{P}(W_j > R_j) \le \epsilon_j$. \status{not checked: drafted
2026-11-15}
) in Section 3. Uses the
% assumption, definition, theorem, corollary and remark environments and the
% \status macro defined there. Citation keys are those of paper_a.bib.
%
% Proofs other than Theorem 1's are in paper_a_proofs.tex (the appendix).
% Results held until probe P1 decides the chemotherapy design (joint radius,
% certified level, simulated level) are in paper_a_guarantee_held.tex, which
% is not \input.
%
% \status tags. JS checked by hand the version committed in 40acea5: Theorem 1,
% Corollaries 1-3 and the closing remark, for Sigma = sigma^2 I. The general
% Sigma, the remark on the choice of norm and everything from "The simulated
% radius" on came later and are not checked.

\begin{outline}
Tags: JS checked by hand the $\Sigma = \sigma^2 I$ version of
Theorem~\ref{thm:selection}, Corollaries~\ref{cor:finite}--\ref{cor:multi}
and the closing remark (commit \texttt{40acea5}). The general $\Sigma$, the
remark on the choice of norm and everything after the closing remark were
added later and are not checked.
\end{outline}

\paragraph{Setting.}
We observe training data $(x_i, y_i)_{i=1}^n$ with $y = y^\star + \varepsilon$,
where $y^\star_i = h^\star(x_i)$ are the noise-free labels. The learned
constraint is $f(x;\theta(y)) \le b$ for $x \in \mathcal{X}$, where $\theta(y)$
denotes the model fitted on labels $y$.

\begin{assumption}[Label noise]\label{ass:noise}
\status{checked by JS for $\Sigma = \sigma^2 I$; general $\Sigma$ not checked}
The noise is $\varepsilon = \Sigma^{1/2}\xi$, where $\Sigma \succeq 0$ is known
and $\xi_1,\dots,\xi_n$ are independent, zero-mean and sub-Gaussian with
variance proxy $1$. We write $\|v\|_\Sigma = \|\Sigma^{1/2}v\|_2$.
\end{assumption}

With $\Sigma = \sigma^2 I$ the noise terms are independent with variance proxy
$\sigma^2$ and $\|v\|_\Sigma = \sigma\|v\|_2$. A diagonal $\Sigma$ allows
unequal noise levels, and a full one correlated noise.

\begin{assumption}[Affine predictor]\label{ass:affine}
For all $x \in \mathcal{X}$, $f(x;\theta(y)) = s(x)^\top y + c(x)$, where
$s(x) \in \mathbb{R}^n$ and $c(x)$ do not depend on $\varepsilon$.
\end{assumption}

\begin{definition}[Clean-label model]\label{def:clean}
Let $\hat\theta = \theta(y)$ and $\theta^\star = \theta(y^\star)$. The
clean-label model is $f(\cdot\,;\theta^\star)$.
\end{definition}

Since $s(x)$ and $c(x)$ do not depend on $\varepsilon$,
Assumption~\ref{ass:affine} applies to $y^\star = y - \varepsilon$ as well as to
$y$, and subtracting the two gives, for all $x \in \mathcal{X}$,
\begin{equation}\label{eq:identity}
  f(x;\theta^\star) = f(x;\hat\theta) - s(x)^\top \varepsilon .
\end{equation}

For a radius $R > 0$ fixed independently of $\varepsilon$, the robust
counterpart of the learned constraint is
\begin{equation}\label{eq:rc}
  f(x;\hat\theta) + R\,\|s(x)\|_\Sigma \le b .
\end{equation}
The radius $R$ is in units of the noise's scale: with $\Sigma = \sigma^2 I$ the
margin is $R\sigma\|s(x)\|_2$.

\begin{theorem}\label{thm:selection}
\status{checked by JS for $\Sigma = \sigma^2 I$; general $\Sigma$ not checked}
Suppose Assumptions~\ref{ass:noise} and~\ref{ass:affine} hold, and let $x^\star$
be any solution satisfying \eqref{eq:rc}, possibly chosen as a function of the
data. Then
\[
  \mathbb{P}\bigl(f(x^\star;\theta^\star) > b\bigr) \le \mathbb{P}(W > R),
  \qquad
  W = \sup_{x \in \mathcal{X},\, \|s(x)\|_\Sigma \neq 0}
      \frac{-s(x)^\top \varepsilon}{\|s(x)\|_\Sigma} .
\]
\end{theorem}

\begin{proof}
Suppose $f(x^\star;\theta^\star) > b$. By \eqref{eq:identity} and
\eqref{eq:rc},
\[
  -s(x^\star)^\top \varepsilon
  = f(x^\star;\theta^\star) - f(x^\star;\hat\theta)
  > b - f(x^\star;\hat\theta)
  \ge R\,\|s(x^\star)\|_\Sigma .
\]
The left side equals $-(\Sigma^{1/2}s(x^\star))^\top\xi$ and is nonzero, so
$\|s(x^\star)\|_\Sigma \neq 0$ and $W > R$.
\end{proof}

\begin{remark}[Choice of norm]\label{rem:norm}
\status{not checked: added after JS's check}
The proof uses only the inequality
$-s(x^\star)^\top\varepsilon > R\,\|s(x^\star)\|_\Sigma$, so
Theorem~\ref{thm:selection} holds verbatim for any norm $N$ on $\mathbb{R}^n$
in place of $\|\cdot\|_\Sigma$, both in \eqref{eq:rc} and in $W$. The norm
enters only through the distribution of $W$. Under $\|\cdot\|_\Sigma$ the
standardized error at $x$ is
$-s(x)^\top\varepsilon/\|s(x)\|_\Sigma = -u(x)^\top\xi$ with $u(x)$ a unit
vector, so it has the same tail at every $x$; Corollaries~\ref{cor:finite}
and~\ref{cor:subspace} use exactly this. Under another norm the standardized
error at $x$ is that one multiplied by
$\rho(x) = \|s(x)\|_\Sigma / N(s(x))$. Bounding the corresponding $W_N$
through $W$, as $W_N \le \bar\rho\max(W,0)$ with
$\bar\rho = \sup_{x\in\mathcal{X}} \rho(x)$, turns a radius $R$ from these
corollaries into $\bar\rho R$, and the margin
$\bar\rho R\,N(s(x)) \ge R\,\|s(x)\|_\Sigma$ is at least as large at every
$x$. For example, noise bounded as $|\varepsilon_i| \le c$ gives the box
counterpart $f(x;\hat\theta) + c\,\|s(x)\|_1 \le b$, which holds with
probability one; but $\|s(x)\|_1$ can be up to $\sqrt{n}$ times
$\|s(x)\|_2$.
\end{remark}

\begin{corollary}[Tree ensembles]\label{cor:finite}
\status{checked by JS for $\Sigma = \sigma^2 I$; general $\Sigma$ not checked}
If $s(x)$ takes at most $K$ distinct values on $\mathcal{X}$, then
$R = \sqrt{2\ln(K/\epsilon)}$ guarantees
$\mathbb{P}\bigl(f(x^\star;\theta^\star) > b\bigr) \le \epsilon$.
\end{corollary}

\begin{corollary}[Linear models]\label{cor:subspace}
\status{checked by JS for $\Sigma = \sigma^2 I$; general $\Sigma$ not checked}
Suppose $s(x)$ lies in a subspace $S \subseteq \mathbb{R}^n$ of dimension $d$
for all $x \in \mathcal{X}$, and $\xi \sim N(0, I)$, that is,
$\varepsilon \sim N(0,\Sigma)$. Then $R = \sqrt{\chi^2_{d,1-\epsilon}}$
guarantees $\mathbb{P}\bigl(f(x^\star;\theta^\star) > b\bigr) \le \epsilon$,
where $\chi^2_{d,1-\epsilon}$ is the $(1-\epsilon)$-quantile of the chi-square
distribution with $d$ degrees of freedom.
\end{corollary}

\begin{corollary}[Multiple learned constraints]\label{cor:multi}
\status{checked by JS for $\Sigma = \sigma^2 I$; general $\Sigma$ not checked}
Consider $m$ learned constraints, each satisfying
Assumptions~\ref{ass:noise} and~\ref{ass:affine} with its own $s_j$,
$\Sigma_j$ and radius $R_j$. If each $R_j$ is chosen as in
Corollary~\ref{cor:finite} or~\ref{cor:subspace} at level $\epsilon/m$, then
the probability that $x^\star$ violates at least one clean-label constraint is
at most $\epsilon$. No assumption is required on the dependence of the noise
across constraints.
\end{corollary}

\begin{remark}
\status{checked by JS for $\Sigma = \sigma^2 I$}
(i) For ridge regression, $s(x) = Z(Z^\top Z + \lambda I)^{-1}\phi(x)$, so $S$
is the column space of the training matrix $Z$ and $d = \mathrm{rank}(Z)$.
(ii) A lower-bound constraint $f(x;\theta) \ge b$ is handled symmetrically,
with robust counterpart $f(x;\hat\theta) - R\|s(x)\|_\Sigma \ge b$ and the same
radii.
(iii) The guarantee concerns the propagation of label noise only; the
approximation error $h^\star - f(\cdot\,;\theta^\star)$ is not covered.
(iv) For models whose structure is selected from data (active sets, splits,
hidden layers), Assumption~\ref{ass:affine} requires the structure to be
chosen independently of $\varepsilon$, for instance by sample splitting.
\end{remark}

\paragraph{The simulated radius.}
Corollaries~\ref{cor:finite} and~\ref{cor:subspace} bound $W$ without using
$\mathcal{X}$: they protect against every direction among the vectors
$\Sigma^{1/2}s(x)$, whether or not a decision in $\mathcal{X}$ reaches it. When
the distribution of $\xi$ is known, the radius can instead be read off
simulated copies of $W$, and an order statistic of the copies is valid with no
sampling-error term. Write $W = w(\xi)$, where
\[
  w(z) = \sup_{x \in \mathcal{X},\, \|s(x)\|_\Sigma \neq 0}
    \frac{-(\Sigma^{1/2}s(x))^\top z}{\|s(x)\|_\Sigma},
  \qquad z \in \mathbb{R}^n ,
\]
assumed measurable. Under Assumption~\ref{ass:affine}, $s(\cdot)$ does not
depend on $\varepsilon$, so $w$ is a fixed function, known once the model's
structure is. Evaluating $w(z)$ is a maximization over $\mathcal{X}$, a
mixed-integer program when $\mathcal{X}$ is mixed-integer.

\begin{proposition}[Simulated radius]\label{prop:simulated}
\status{not checked: drafted 2026-11-15}
Suppose Assumptions~\ref{ass:noise} and~\ref{ass:affine} hold and the
distribution of $\xi$ is known. Draw $\xi^{(1)},\dots,\xi^{(N)}$ from it,
independently of each other and of $\varepsilon$, and let
$\overline W_1,\dots,\overline W_N$ be any numbers with
$\overline W_i \ge w(\xi^{(i)})$. For $\epsilon \in (0,1)$ and
$N \ge 1/\epsilon - 1$, let $k = \lceil (N+1)(1-\epsilon) \rceil$ and let
$\hat R$ be the $k$-th smallest of $\overline W_1,\dots,\overline W_N$. Then
any $x^\star$ satisfying \eqref{eq:rc} at $R = \hat R$, possibly chosen as a
function of the data and of the simulation, satisfies
$\mathbb{P}\bigl(f(x^\star;\theta^\star) > b\bigr) \le \epsilon$, where the
probability is over the noise and the simulation jointly.
\end{proposition}

\begin{remark}[Computing the radius]\label{rem:simulated}
\status{not checked: drafted 2026-11-15}
(i) Each $\overline W_i$ must be an upper bound on $w(\xi^{(i)})$. The value
of a feasible point of the maximization, or the largest standardized error
over a sample of decisions, is a lower bound, and an order statistic of lower
bounds can fall below the $(1-\epsilon)$-quantile of $W$, which makes the
radius too small. A solver's dual bound is an upper bound at any time limit,
so stopping a maximization early makes the radius larger, never invalid.
(ii) The proposition needs the distribution of $\xi$, not only a variance
proxy: the argument requires the simulated copies to have the law of the
true $W$. The experiments use $\xi \sim N(0, I)$. An estimated $\Sigma$ is
not covered yet. \tbd{T2}
\end{remark}

Corollary~\ref{cor:multi} holds with unequal levels $\epsilon_j$ summing to
$\epsilon$, and with any radius valid at level $\epsilon_j$ for constraint $j$
alone, including the simulated radius of
Proposition~\ref{prop:simulated}: its proof uses only
$\mathbb{P}(W_j > R_j) \le \epsilon_j$. \status{not checked: drafted
2026-11-15}


\begin{outline}
\begin{itemize}
  \item \textbf{Simulated radius} (the radius the experiments use). Draw $N$
    noise vectors, compute $W$ for each by maximizing over $\mathcal{X}$, and
    take the $\lceil (N+1)(1-\epsilon) \rceil$-th smallest. The true $W$ and
    the $N$ draws are exchangeable, so the violation probability is at most
    $\epsilon$. Each maximization is solved with the solver's upper bound.
    Smaller than the closed-form radii because the feasible decisions reach
    only part of the model's parameter space. \tbd{to prove: short}
    Precedent to credit here: for a linear model over a box this is the
    simulation-based band of \citet{liu2005simulation}, whose per-draw
    supremum is a quadratic program; ours is a global MIP over the decision
    set, and must use the solver's bound, not its incumbent, or the radius
    is anticonservative.
  \item \textbf{Positioning, one paragraph.} The proof of
    Theorem~\ref{thm:selection} is the projection argument of post-selection
    inference \citep{berk2013valid, andrews2024inference}: $W$ is the maximum
    of standardized errors over the directions $s(x)/\|s(x)\|_2$,
    $x \in \mathcal{X}$, and Corollary~\ref{cor:subspace} is Scheff\'e's
    radius \citep{scheffe1953method}. Uniform bounds of this kind also
    underlie safe Bayesian optimization and learning-based control;
    \citet{fiedler2021practical} is the closest, and targets the true
    function, so its radius needs a bound on the function's norm that ours
    does not.
  \item One table of radii: pointwise (not valid at $x^\star$), closed-form,
    simulated. \tbd{first numbers on the reactor at $\epsilon = 0.1$: cubic
    ridge 1.28$\sigma$, 8.4$\sigma$, about 3.7$\sigma$; network (last layer,
    24 dimensions) 1.28$\sigma$, 5.8$\sigma$, about 3.5$\sigma$}. The
    simulated radius must use the solver's bound on each $W$: enumeration over
    sampled points understated it by 0.10 to 0.24$\sigma$ on the network in
    every one of ten draws.
  \item Several constraints: the simulated joint radius. \tbd{to prove} With
    fixed context variables, $W$'s supremum is over the treatment slices of
    the contexts prescribed for, not the whole box.
  \item \textbf{Certified level of any decision.} \tbd{to write; the argument
    is Theorem~\ref{thm:selection}'s proof} With
    $\hat r(x) = \min_j (b_j - f_j(x;\hat\theta_j)) / \|\Sigma_j^{1/2}
    s_j(x)\|_2$, violation implies $W > \hat r(x^\star)$, so
    $\mathbb{P}(\text{violation and } \hat r(x^\star) \ge r) \le
    \mathbb{P}(W > r)$ for every $r$, however $x^\star$ was chosen. This
    certifies any method's decision on one scale, and it is the fallback when
    the RC at the target radius is infeasible. The statement is joint, not
    conditional on reaching $\hat r$.
  \item Estimated $\sigma$: guarantee with an upper confidence bound.
    \tbd{to prove}
  \item Learned structure: what honest fitting means for trees
    \citep{wager2018estimation} and for the hidden layers of a network; the
    method runs on any frozen structure.
  \item Remark on the true function: extending the guarantee from
    $f(\cdot\,;\theta^\star)$ to $h^\star$ needs an assumption on $h^\star$;
    Section~\ref{sec:experiments} measures the gap.
\end{itemize}
\end{outline}

% ---------------------------------------------------------------------------
\section{Solving the robust counterpart}\label{sec:solving}

\begin{outline}
\begin{itemize}
  \item Linear models and last-layer networks: one second-order cone of the
    dimension of the feature vector; hidden layers by the standard ReLU
    embedding. \tbd{first tests on the reactor: cubic ridge solves in under
    one second; a (32, 32) network in 4 to 10 seconds on average; one
    solver-certified draw of $W$ takes 4 to 19 seconds on the network and
    about 43 on cubic ridge}
  \item Trees and forests: $\|s(x)\|_2$ through the leaf binaries; directed
    cuts \citep{mutapcic2009cutting}; at most $K$ iterations. \tbd{to prove:
    short} Reformulation against cuts compared as in
    \citet{bertsimas2016reformulation}.
  \item Certifying any decision: one evaluation of $\hat r(x)$ gives the
    level at which a decision is certified, whichever method produced $x$.
  \item When the RC is infeasible: maximize $t$ subject to the RC at radius
    $t$, by bisection over feasibility problems, and report the level
    $\mathbb{P}(W > \hat r)$; with no useful $t$, report abstention.
\end{itemize}
\end{outline}

% ---------------------------------------------------------------------------
\section{Baselines}\label{sec:baselines}

\begin{outline}
\begin{itemize}
  \item Nominal; constant margin; wrapper \citep{maragno2025mixed}; C-MICL
    \citep{ovalle2025cmicl} at its published settings.
  \item One model per problem for every method.
  \item Each method has one conservatism parameter. Two comparisons: as
    prescribed for a 90\% target, and at equal ground-truth feasibility.
\end{itemize}
\end{outline}

% ---------------------------------------------------------------------------
\section{Computational experiments}\label{sec:experiments}

\begin{outline}
Nothing in this section has been run on the paper's final design. The lead
model on the reactor is not yet decided (see Next steps). Preliminary probes
are summarized after the plan below.
\begin{itemize}
  \item \textbf{Membrane reactor design.} One learned constraint, ODE oracle,
    $\sigma = 2$; the cost vectors of \citet{ovalle2025cmicl} with the noise
    redrawn for each. Table 1: feasibility, cost, fraction solved, solve time,
    rows used, per method at target 0.9. Figure 2: cost against ground-truth
    feasibility.
  \item \textbf{Is the guarantee tight enough to use?} Same runs at three
    radii and three values of $\epsilon$. Figure 3: violation rate against
    $\epsilon$, against the clean-label model and the ODE.
  \item \textbf{Chemotherapy regimen design.} Real trial data, five learned
    toxicity constraints, the protocol of \citet{maragno2025mixed}; $\sigma$
    per arm from the binomial sampling error. Where the RC is infeasible, the
    certified fallback. Table 2: fraction of cohorts certified at
    $\epsilon \in \{0.05, 0.1, 0.2\}$ for every method, fraction prescribed,
    toxicity limits met under the judge, predicted survival.
  \item \textbf{World Food Programme diet.} Table 3, same columns as Table 1.
\end{itemize}
\end{outline}

\begin{outline}
\textbf{Preliminary results (probes, 2026-11-14; not the paper's design).}
Reactor, the 100 cost vectors of \citet{ovalle2025cmicl}, $\sigma = 2$ known,
noise redrawn for every run, $\epsilon = 0.1$. Feasibility against the
clean-label model / against the ODE; standard errors at most 0.05.
\texttt{experiments/probe\_rc\_ridge.py}, \texttt{results/paper\_a/prelim/}.

\smallskip
\begin{tabular}{@{}lcccc@{}}
\toprule
Model & Nominal & Pointwise & Simulated & Closed-form \\
\midrule
Cubic ridge                    & 0.49 / 0.66 & 0.90 / 0.80 & 1.00 / 0.95 & 1.00 / 1.00 \\
Network, hidden layers frozen  & 0.60 / 0.18 & 0.91 / 0.22 & 1.00 / 0.52 & 1.00 / 0.79 \\
Network, sample split          & 0.57 / 0.25 & 0.92 / 0.28 & 1.00 / 0.42 & 1.00 / 0.61 \\
Network, standard fitting      & 0.32 / 0.23 & 0.42 / 0.29 & 0.76 / 0.33 & 0.98 / 0.45 \\
\bottomrule
\end{tabular}
\smallskip

\begin{itemize}
  \item The RC delivers on cubic ridge at 4.2\% of the objective range over
    nominal (closed-form: 8.6\%).
  \item On the network it does not, and the cause is bias, not label noise:
    at the nominal optimum the clean-label network overstates the ODE in 80
    to 90\% of runs, by 1.6 to 2.8 flow units on average; cubic ridge's error
    there has mean zero. The optimizer selects on misfit, which a label-noise
    margin cannot see.
  \item Standard fitting breaks the guarantee itself: 0.76 against the
    clean-label model at a 0.90 target.
  \item No selection effect against the clean-label model on cubic ridge: at
    $\epsilon \in \{0.05, 0.1, 0.2\}$ the pointwise radius is on target (0.95,
    0.90, 0.77). The simulated radius never violates the clean-label model,
    and against the ODE holds 0.93 to 0.97 where pointwise holds 0.76 to
    0.81: its extra margin pays for the spread of the bias.
\end{itemize}

\textbf{C-MICL by surrogate family} (their fixed settings, their 100 cost
vectors, 4 to 5 noise redraws; \texttt{experiments/probe\_cmicl\_family\_coverage.py}).
Truth feasible at a random point inside C-MICL's feasible region / at the
optimizer's $x^\star$; and coverage at truly feasible points inside the region
/ anywhere.

\smallskip
\begin{tabular}{@{}lcc@{}}
\toprule
Surrogate & Feasible: random point / $x^\star$ & Coverage: in region / anywhere \\
\midrule
MLP (32, 32)           & 0.99 / 0.60 & 0.93 / 0.97 \\
Gradient boosting      & 0.96 / 0.72 & 0.66 / 0.94 \\
Random forest          & 0.93 / 0.59 & 0.49 / 0.93 \\
Single tree, depth 5   & 0.45 / 0.45 & 0.66 / 0.96 \\
\bottomrule
\end{tabular}
\smallskip

\begin{itemize}
  \item Their Assumption 4.1 fails for every tree family (coverage drops
    inside the feasible region) and nearly holds for the MLP.
  \item Their guarantee nonetheless holds at random feasible points for the
    MLP and both ensembles; at the optimizer's choice it fails for all four.
    The shortfall is selection, which no exchangeability argument covers.
  \item C-MICL's feasible region is 0.2 to 2.8\% of the design domain and
    holds 0 to 12 of its 200 calibration rows: the quantile is calibrated
    where no decision can go.
  \item Nominal tree decisions sit at the cheapest point of their leaf cell
    in every run, and most of their optimistic error (3.9 to 11.4 flow units)
    is position within the cell. A tree RC, whose margin is constant on a
    cell, cannot see this either.
\end{itemize}
\end{outline}

% ---------------------------------------------------------------------------
\section{Discussion}

\begin{outline}
\begin{itemize}
  \item When to use which safeguard.
  \item Limitations: the guarantee is against the clean-label model; learned
    structure needs honest fitting for the theorem; models not affine in the
    labels.
  \item Classification, one paragraph: linear probability models and
    classification forests are affine in Bernoulli labels, whose noise is
    bounded; logistic models are not. Future work.
  \item Future work: feature noise, model selection for embedding, recourse.
\end{itemize}
\end{outline}

% ---------------------------------------------------------------------------
% INTERNAL: for the authors only. Remove before the draft is circulated.
\section*{Next steps (internal draft)}
\addcontentsline{toc}{section}{Next steps (internal draft)}

\begin{tcolorbox}[breakable, colback=blue!3, colframe=blue!40!black,
  title={For JS and DB only; remove before circulating},
  fonttitle=\small\bfseries, left=4pt, right=4pt, top=3pt, bottom=3pt]

\textbf{Decisions for DB.}
\begin{enumerate}
  \item \textbf{Lead model on the reactor.} Cubic ridge works (nominal 0.66,
    RC 0.95 against the ODE), but no reference paper embeds it. The network of
    \citet{ovalle2025cmicl} fails because of its bias (RC 0.42 to 0.52). The
    third option is their LinearDT surrogate (a regression in each leaf): it
    is at the noise floor on this instance, so its bias should be small, but
    it is untested and needs a MIP embedding. Recommendation: cubic ridge as
    lead, LinearDT probed next, the network as the example of where a
    label-noise margin is not enough.
  \item \textbf{Bias: a remark or part of the method?} Two results now show
    the optimizer selecting on model bias, not label noise: the network (bias
    optimistic in 80 to 90\% of runs at $x^\star$) and trees (within-cell
    error of 4 to 11 flow units at the nominal $x^\star$). Either the paper
    restricts its claim to low-bias models and says so, or the bias margin
    $B(x)$ (T5) becomes part of the method.
  \item \textbf{Figure 3's argument.} On cubic ridge the pointwise radius is
    already on target against the clean-label model at every $\epsilon$, so
    the probe shows no price of selection. Either find the setting where
    selection bites (a cost nearly flat along the boundary, smaller $n$,
    higher dimension) or reframe the figure around robustness to bias.
  \item \textbf{How much of the C-MICL analysis goes in.} Assumption 4.1
    fails for trees; the guarantee fails at $x^\star$ for every surrogate.
    Options: a paragraph in Section~\ref{sec:baselines}, or a short
    subsection. It is the clearest evidence for the paper's premise.
  \item \textbf{Several constraints against a ground truth.} No experiment
    tests the joint radius against a truth (the reactor has one learned
    constraint, chemotherapy has no truth). Accept and say so, or add a
    second ODE output as a constraint on the reactor.
  \item \textbf{Random-bank cutting planes as a baseline?} Recommendation:
    no. No guarantee, two oracle-tuned parameters, hours per cell.
\end{enumerate}

\textbf{Experiments, in order.}
\begin{enumerate}
  \item Probe P1 on chemotherapy: ridge on the training arms, the max-$t$
    fallback per test cohort, binomial $\sigma$ against residual sd. Go/no-go
    for H3 before the forest RC (I6), the chemo runner (I7) and the certified
    radius (I9) are built.
  \item Once the lead model is fixed: E1 on the reactor with the
    solver-certified simulated radius, all baselines on the shared model,
    plus C-MICL fitted as published. 100 redraws; 300 for the Table 1 rows if
    the RC and C-MICL land within 0.05.
  \item LinearDT: embedding, then the same probe as the network (job
    29141177's design).
  \item A selection check for Figure 3: how far $x^\star$ moves across
    redraws, and one cost direction nearly flat along the boundary.
  \item Minor: the C-MICL family probe crashes when no sampled point is
    C-MICL-feasible (two of 20 runs); guard it before any rerun.
\end{enumerate}

\textbf{Theory, in order.} T1 (simulated radius, single and joint); T6
(certified level of any decision; needed for the chemo fallback); T2
(estimated and per-arm $\sigma$, including the binomial plug-in); T3 (cuts
terminate). Check whether C-MICL's coverage event is for the noisy label or
the true value (their App.~B): with the true value, the infeasible half of
Assumption 4.1 can hold only if C-MICL is never wrong.

\textbf{Writing.} Sections~\ref{sec:method} and~\ref{sec:guarantee} are prose.
The introduction's numbers wait on the lead model; Section~\ref{sec:solving}
and the experiments wait on E1 and P1.
\end{tcolorbox}

\bibliographystyle{plainnat}
\bibliography{paper_a}

\appendix
\section{Proofs and formulations}
\begin{outline}
Sub-Gaussian version of the subspace radius; formulations per model class; the
bias margin for Lipschitz $h^\star$ (if proved); experimental detail and three
checks on the reactor (standard against honest forest, non-Gaussian noise,
cubic ridge).
\end{outline}

\end{document}
